\begin{frame}[t, fragile]{Types classes in Haskell.}

Haskell has a number of different type classes. Some examples of the different type classes which are available are \mintinline[fontsize=\normalsize]{haskell}{Num},
\mintinline[fontsize=\normalsize]{haskell}{Show},
\mintinline[fontsize=\normalsize]{haskell}{Read},
\mintinline[fontsize=\normalsize]{haskell}{Eq}, and
\mintinline[fontsize=\normalsize]{haskell}{Ord}. \\

\textbullet{} Creating a new type that is an instance of a type class. \\

Imagine that we want to create a new type called \mintinline[fontsize=\normalsize]{haskell}{Person}, and that we define
it as follows;

\begin{minted}[linenos=false]{haskell}
data Person = Person String Int
\end{minted}

This says that an instance of the \mintinline[fontsize=\normalsize]{haskell}{Person} type will contain a
\mintinline[fontsize=\normalsize]{haskell}{String} called name and an \mintinline[fontsize=\normalsize]{haskell}{Int}
called age. The trouble is, if we try and use the \mintinline[fontsize=\normalsize]{haskell}{print} function on an
instance of this type, then the compiler will complain that the type \mintinline[fontsize=\normalsize]{haskell}{Person}
is not an instance of the type class \mintinline[fontsize=\normalsize]{haskell}{Show}.

To remedy this, our type \mintinline[fontsize=\normalsize]{haskell}{Person} must implement the type class
\mintinline[fontsize=\normalsize]{haskell}{Show} as follows;

\begin{minted}[linenos=true]{haskell}
data Person = Person String Int
    deriving Show
\end{minted}

\end{frame}


% New slide.

\begin{frame}[t, fragile]{Types classes in Haskell.}

Note line 2 of the previous code snippet; in particular, the use of the keyword \mintinline[fontsize=\normalsize]{haskell}{deriving}.
This keyword instructs the Haskell compiler to generate an instance of one or more type classes for this type, rather
than you having to write the instance yourself.

After line 1, \mintinline[fontsize=\normalsize]{haskell}{Person} is just a new type. It is not yet an instance of the
type class \mintinline[fontsize=\normalsize]{haskell}{Show}. \\

\textbullet{} Deriving (from) multiple type classes. \\

You can derive several type classes at once using code which is similar to the following;

\begin{minted}[linenos=true]{haskell}
data Person = Person String Int
    deriving (Show, Eq)
\end{minted}

This is vaguely like the concept of multiple inheritance.
\end{frame}
